% Part: first-order-logic
% Chapter: syntax-and-semantics
% Section: subformulas

\documentclass[../../../include/open-logic-section]{subfiles}

\begin{document}

\olfileid{fol}{syn}{sbf}

\olsection{\printtoken{P}{subformula}}

\begin{explain}
ډېر ځله د هغو !!{formula}s په باب خبرې کول ګټور وي چې يو راکړل شوے
!!{formula} ``جوړوي''۔ دې ته د هغه \emph{!!{subformula}s} وايو۔ هر
!!{formula}~$!A$ د خپل ځان !!a{subformula} هم ګڼل کېږي؛ د $!A$ هر هغه
فرعي فارمول چې پخپله $!A$ نۀ وي، \emph{حقيقي !!{subformula}} دے۔
\end{explain}

\begin{defn}[سمدستي !!^{subformula}]
که $!A$ !!a{formula} وي، د $!A$ \emph{سمدستي !!{subformula}s} په
استقرايي ډول داسې تعريفېږي:
\begin{enumerate}
\item اټومي !!{formula}s سمدستي !!{subformula}s نۀ لري۔

\tagitem{prvNot}{\indcase{!A}{\lnot !B}{د $\indfrm$ يوازينے سمدستي
    !!{subformula} همدا~$!B$ دے۔}}{}

\item \indcase{!A}{(!B \ast !C)}{د $\indfrm$ سمدستي
  !!{subformula}s همدا $!B$ او $!C$ دي ($\ast$ له دوه ځايه
  نښلوونکو څخه هر يو کېدے شي)۔}

\tagitem{prvAll}{\indcase{!A}{\lforall[x][!B]}{د $\indfrm$ يوازينے
    سمدستي !!{subformula} همدا~$!B$ دے۔}}{}

\tagitem{prvEx}{\indcase{!A}{\lexists[x][!B]}{د $\indfrm$ يوازينے
    سمدستي !!{subformula} همدا~$!B$ دے۔}}{}
\end{enumerate}
\end{defn}

\begin{defn}[حقيقي !!^{subformula}]
که $!A$ !!a{formula} وي، د $!A$ \emph{حقيقي !!{subformula}s} په
بازګشتي ډول داسې تعريفېږي:
\begin{enumerate}
\item اټومي !!{formula}s حقيقي !!{subformula}s نۀ لري۔

\tagitem{prvNot}{\indcase{!A}{\lnot !B}{د $\indfrm$ حقيقي
    !!{subformula}s همدا~$!B$ او د~$!B$ ټول حقيقي !!{subformula}s دي۔}}{}

\item \indcase{!A}{(!B \ast !C)}{د $\indfrm$ حقيقي
  !!{subformula}s همدا $!B$، $!C$، د $!B$ ټول حقيقي !!{subformula}s
  او د~$!C$ ټول حقيقي فرعي فارمولونه دي۔}

\tagitem{prvAll}{\indcase{!A}{\lforall[x][!B]}{د $\indfrm$ حقيقي
    !!{subformula}s همدا~$!B$ او د~$!B$ ټول حقيقي
    !!{subformula}s دي۔}}{}

\tagitem{prvEx}{\indcase{!A}{\lexists[x][!B]}{د $\indfrm$ حقيقي
    !!{subformula}s همدا~$!B$ او د~$!B$ ټول حقيقي
    !!{subformula}s دي۔}}{}
\end{enumerate}
\end{defn}

\begin{defn}[!!^{subformula}]
د $!A$ !!{subformula}s پخپله $!A$ او د هغه ټول حقيقي
!!{subformula}s دي۔
\end{defn}

\begin{explain}
پام وکړئ چې سمدستي !!{subformula}s او حقيقي !!{subformula}s مو په
لږ توپير تعريف کړي دي۔ په لومړي حالت کښې مو د~$!A$ د هرې ممکنې بڼې
دپاره د هغه سمدستي !!{subformula}s نېغ په نېغه تعريف کړي دي۔ دا د
حالتونو له مخې يو څرګند تعريف دے، او حالتونه ئې د !!{formula}s د سټ
له استقرايي تعريف سره سمون لري۔ په دويم حالت کښې مو هم د ټولو
!!{formula}s د سټ د تعريف طريقه تعقيب کړې ده، خو په هر حالت کښې مو د
کوچنيو !!{formula}s $!B$ او $!C$ د حقيقي !!{subformula}s تر څنګ،
پخپله دا !!{formula}s هم ور شامل کړي دي۔ په دې سره تعريف \emph{بازګشتي}
کېږي۔ په عمومي ډول، پر يوه استقرايي تعريف شوي سټ (زموږ په حالت کښې
!!{formula}s) د تابع تعريف هله بازګشتي وي چې د تابع د تعريف حالتونه
پخپله هماغه تابع وکاروي۔ خو د ښه تعريف دپاره بايد ډاډه شو چې د تابع
قيمتونه تل يوازې د هغو مدخلو دپاره کاروو چې تر تعريفېدونکې مدخل
``مخکښې'' راځي---زموږ په حالت کښې، د $(!B \ast !C)$ د ``حقيقي
!!{subformula}'' د تعريف پر وخت يوازې د ``مخکښو'' !!{formula}s $!B$
او $!C$ حقيقي !!{subformula}s کاروو۔
\end{explain}

\begin{prop}
\ollabel{prop:subfrm-trans}
فرض کړئ $!B$ د $!A$ يو فرعي فارمول او $!C$ د $!B$ يو فرعي فارمول
دے۔ بيا $!C$ د $!A$ يو فرعي فارمول دے۔ په بله وينا، د فرعي فارمول
اړيکه انتقالي ده۔
\end{prop}

\begin{prob}
\olref[fol][syn][sbf]{prop:subfrm-trans} ثابته کړئ۔
\end{prob}

\begin{prop}
\ollabel{prop:count-subfrms}
فرض کړئ $!A$ داسې فارمول دے چې $n$ نښلوونکي او سورونه لري۔ بيا $!A$
تر ډېره $2n+1$ فرعي فارمولونه لري۔
\end{prop}

\begin{prob}
\olref[fol][syn][sbf]{prop:count-subfrms} ثابته کړئ۔
\end{prob}

\end{document}
